% Generated by scripts/build-paper.py — edit paper/src/survey.src.md
\section{Review method}\label{sec:review-method}

\subsection{Search}

The snapshot of 23 September 2026 ran 17 accepted arXiv API queries. They returned 225 hits and 197 unique records after de-duplication by unversioned arXiv ID (Appendix A). Queries combined the brand and interface vocabulary (\texttt{jev}, \texttt{typesafe}, “System One”, “typed decision”, “calibrated decisions”) with open-model names, survey terms and a date-bounded “decision model” query. The last matters: arXiv \texttt{all:} searches metadata rather than full text, and it found a core study that brand-name queries missed \citep{arxiv2609_24574}. One compound query reported 1,232,608 results and was stopped by rate limiting; it was rejected rather than screened and replaced by narrow title/abstract queries. Six background references came from these queries and 38 more were retrieved by explicit ID as a purposive seed for theory and baselines.

A same-day increment re-ran all 17 queries. The totals were identical. It then added 19 expansion queries on interface vocabulary and open-model names; four were rejected by rule because the API parsed their phrases so loosely that they returned thousands of records. It screened 128 new records by title and abstract and added 9 as background references. No new core study appeared, which is expected: arXiv announces new submissions around 00:00 UTC. For repositories, five GitHub searches produced 1,082 unique discovery results. From these we audited 81 catalogues and research guides (79 READMEs at pinned commits) and 30 priority implementation and evaluation resources. The 4,666 outgoing links found in catalogue READMEs are kept as unverified candidates and never counted as projects.

\subsection{Eligibility}

A study is \textbf{core} if it evaluates TypeSafe Jev, a clearly identified Jev-like typed-decision implementation, or a system whose contribution materially depends on such decisions. \textbf{Peripheral} studies fall within reach of the scope but rest on contested, self-reported claims; they stay visible but are not pooled. \textbf{Background} references are adjacent methods with a stated role. We excluded name matches (Japanese encephalitis, JEPA, unrelated uses of “TypeSafe”, other expansions of RLCD) and generic System 1/2 work without a typed-decision link.

\subsection{Reading depth and extraction}

We read the core and peripheral studies in full, checking key tables and limitation sections. Background references were checked against metadata and abstracts; one reranker paper was spot-checked in full \citep{arxiv2606_22807}. Repository evidence comes from READMEs at pinned commits and 12 source files read in 10 repositories. Nothing was executed and no experiment was re-run.

Each checkable statement became an evidence record. A record stores its subject, text, headline and source locator, and an evidence type: author-reported experiment, vendor documentation, vendor-reported result, community report, or code/README inspection. It also stores the metric, value and baseline; the sample size, model version and hardware, each with a reason when unknown; the \textbf{test level} (model test, system test, hybrid, simulation); and the \textbf{measurement scope} (single request, amortized per question, batch, end to end, simulation, author estimate, vendor claim). Each record links to the findings of §7 as \emph{supporting} or \emph{qualifying} evidence. Openness is recorded as six independent fields per study: code visible, weights available, data available, raw predictions available, recomputable, independently reproduced.

\subsection{Synthesis and its limits}

Tasks, metrics, binning and scopes are heterogeneous, so we pool nothing and rank nothing. Findings are stated at the strength their records allow, with qualifying records listed beside supporting ones. The two edge-orchestration studies share authors and a service path, so they form one study family \citep{arxiv2609_23136,arxiv2609_22753}. The review was carried out by one AI-assisted reviewer and is not a registered or PRISMA-compliant systematic review; §12 lists what this implies.
